\id{МРНТИ 61.01.91}{}

\begin{header}
\swa{Р.О.~Орынбасар, А.Е.~Амангалиева, Л. К.~Тастанова, Г.Ж.~Жакупова, А.К.~Жумабекова}{ИСПОЛЬЗОВАНИЕ ТЕХНОГЕННЫХ ОТХОДОВ В КАЧЕСТВЕ НАПОЛНИТЕЛЕЙ ДЛЯ ПОЛИМЕРНЫХ КОМПОЗИТОВ}

\tsp{1}Р.О.~Орынбасар,
\tsp{1}А.Е.~Амангалиева\envelope,
\tsp{1}Л. К.~Тастанова,
\tsp{1}Г.Ж.~Жакупова,
\tsp{2}А.К.~Жумабекова
\end{header}

\begin{affil}
\tsp{1}Актюбинский региональный университет им. К. Жубанова, Актобе, Казахстан,

\tsp{2}Казахский университет технологии и бизнеса им. К. Кулажанова, Астана, Казахстан

\corrauthor{Корреспондент-автор: Ashlen002@gmail.com}
\end{affil}

В условиях увеличения объемов промышленных отходов остро стоит проблема
их утилизации. Одним из перспективных направлений является использование
техногенных отходов в качестве наполнителей для полимерных
композиционных материалов. В данной обзорной статье рассмотрены
современные подходы (2018-2025 гг.) к переработке золы уноса,
металлургических шлаков, строительного мусора, красного шлама, стеклобоя
и других отходов путем введения их в полимерные матрицы (термопласты,
термореактивные смолы и др.). Приведены сведения о влиянии таких
наполнителей на механические, физико-химические и термические свойства
композитов, выполнен их сравнительный анализ. Отдельное внимание уделено
реализованным технологиям и патентам, демонстрирующим практическое
применение отходов (например, заполнение ПВХ-материалов золой уноса для
труб). Проанализирована экологическая и экономическая эффективность
подобной утилизации отходов: снижения нагрузки на полигоны, замещения
дефицитных минеральных наполнителей и удешевления продукции. Упомянуты
классические работы и приведены ссылки на актуальные исследования
последних лет, подтверждающие рост интереса к ``зеленым'' композитам на
основе вторичного сырья. Полученные сведения демонстрируют, что
применение техногенных отходов в полимерных композитах -- востребованный
и оправданный шаг на пути к ресурсосбережению и экологической
безопасности.

{\bfseries Ключевые слова}: техногенные отходы, наполнители, стеклобой,
красный шлам, композит, рециклинг.

\begin{header}
USE OF TECHNOGENIC WASTE AS FILLERS FOR POLYMER COMPOSITES

\tsp{1}R.O.~Orynbasar,
\tsp{1}A.E.~Amangalieva\envelope,
\tsp{1}L.K.~Tastanova,
\tsp{1}G.Zh.~Zhakupova,
\tsp{2}A.K.~Zhumabekova
\end{header}

\begin{affil}
\tsp{1}K. Zhubanov Aktobe Regional University, Aktobe, Kazakhstan,

\tsp{2}K. Kulazhanov Kazakh University of Technology and Business, Astana, Kazakhstan,

e-mail: Ashlen002@gmail.com
\end{affil}

In conditions of increasing volumes of industrial waste, the problem of
their disposal is acute. One of the promising areas is the use of
man-made waste as fillers for polymer composite materials. This review
article examines modern approaches (2018-2025) to the processing of fly
ash, metallurgical slags, construction debris, red sludge, cullet and
other waste by introducing them into polymer matrices (thermo\-plastics,
thermosetting resins, etc.). Information is provided on the effect of
such fillers on the mechanical, physico-chemical and thermal properties
of composites, and their comparative analysis is performed. Special
attention is paid to implemented technologies and patents demonstrating
the practical use of waste (for example, filling PVC materials with fly
ash for pipes). The environmental and economic efficiency of such waste
disposal is analyzed: reducing the load on landfills, replacing scarce
mineral fillers and reducing the cost of products. Classical works are
mentioned and links to current research in recent years are provided,
confirming the growing interest in ``green'' composites based on
recycled materials. The information obtained demonstrates that the use
of man-made waste in polymer composites is a sought-after and justified
step towards resource conservation and environmental safety.

{\bfseries Keywords}: man-made waste, fillers, cullet, red sludge,
composite, recycling.

\begin{header}
ПОЛИМЕРЛІ КОМПОЗИТТЕР ҮШІН ТОЛТЫРҒЫШ РЕТІНДЕ ТЕХНОГЕНДІК ҚАЛДЫҚТАРДЫ ПАЙДАЛАНУ

\tsp{1}Р.О.~Орынбасар,
\tsp{1}А.Е.~Амангалиева,
\tsp{1}Л.К.~Тастанова,
\tsp{1}Г.Ж.~Жакупова,
\tsp{1}А.К.~Жумабекова
\end{header}

\begin{affil}
\tsp{1}Қ.Жұбанов атындағы Ақтөбе өңірлік университеті, Ақтөбе, Қазақстан,
       
\tsp{2}Қ. Құлажанов атындағы Қазақ технология және бизнес университеті, Астана, Қазақстан,

e-mail: Ashlen002@gmail.com
\end{affil}

Өнеркәсіптік қалдықтар көлемінің ұлғаюы жағдайында оларды кәдеге жарату
мәселесі өткір болып тұр. Перспективалы бағыттардың бірі-техногендік
қалдықтарды полимерлі композициялық материалдар үшін толтырғыш ретінде
пайдалану. Бұл шолу мақаласында полимерлі матрицаларға (термопластика,
термореактивті шайырлар және т.б.) енгізу арқылы тасығыштың күлін,
металлургиялық қождарды, құрылыс қоқыстарын, қызыл шламды, шыны тұнбаны
және басқа да қалдықтарды қайта өңдеудің заманауи тәсілдері (2018-2025
жж.) қарастырылған. Мұндай толтырғыштардың композиттердің механикалық,
физика-химиялық және термиялық қасиеттеріне әсері туралы мәліметтер
келтірілген, сонымен бірге олардың салыстырмалы талдауы жасалды.
Қалдықтардың практикалық қолданылуын көрсететін енгізілген технологиялар
мен патенттерге ерекше назар аударылады (мысалы, ПВХ материалдарын
құбырларға арналған күлмен толтыру). Мұндай қалдықтарды жоюдың
экологиялық және экономикалық тиімділігі талданды: полигондарға
жүктемені азайту, жетіспейтін минералды толтырғыштарды ауыстыру және
өнімді арзандату. Классикалық жұмыстар туралы айтылды және қайталама
шикізатқа негізделген "жасыл" композиттерге деген қызығушылықтың артуын
растайтын соңғы жылдардағы өзекті зерттеулерге сілтемелер келтірілді.
Алынған мәліметтер полимерлі композиттерде техногендік қалдықтарды
қолдану ресурстарды үнемдеу және экологиялық қауіпсіздік жолындағы
сұранысқа ие және негізделген қадам екенін көрсетеді.

{\bfseries Түйін сөздер}: техногендік қалдықтар, толтырғыштар, әйнек, қызыл
шлам, композит, қайта өңдеу.

\begin{multicols}{2}
{\bfseries Введение.} Ускорение промышленного производства неизбежно
сопровождается ростом объёмов техногенных отходов, для которых требуется
технологически и экологически обоснованная утилизация {[}1, 2{]}. К
числу наиболее массовых потоков относятся зола уноса, образующаяся при
сжигании угля: ежегодно формируется свыше 500 млн тонн, причём более 200
млн тонн по-прежнему не находит полезного применения {[}3{]}.

Сопоставимую проблему создаёт красный шлам - побочный продукт получения
глинозёма: его выпуск оценивают примерно в 150 млн тонн в год, и
значительная часть материала длительно накапливается в хвостохранилищах.
Металлургический сектор, в свою очередь, генерирует миллионы тонн
доменных и сталеплавильных шлаков; строительная индустрия - значительные
объёмы строительного мусора; стекольное производство - тонны стеклобоя.

Во многих случаях такие отходы складируются на полигонах, что приводит к
изъятию земельных ресурсов и повышает риск загрязнения окружающей среды
тяжёлыми металлами и другими токсикантами {[}2{]}. По имеющимся оценкам,
при сохранении текущих тенденций без заметного улучшения переработки
общий объём отходов может превысить 3,4 млрд тонн к 2050 году {[}1{]}.
Параллельно увеличивается потребность в материалах, получение которых
связано с расходованием природного сырья и энергозатратами. Поэтому
вовлечение техногенных отходов во вторичный оборот рассматривается не
только как природоохранная мера, но и как способ частично заместить
первичные ресурсы.

Перспективным направлением здесь выступают полимерные композиционные
материалы, в которых дисперсные минеральные фракции (в том числе
полученные из отходов) используются как наполнители. Такие добавки
позволяют уменьшить долю дорогостоящей полимерной матрицы, снизить
себестоимость и одновременно целенаправленно корректировать свойства
композита (прочность, жёсткость, износостойкость, термостойкость и др.)
{[}2{]}. При корректном подборе состава и предварительной обработке
техногенное сырьё может работать как функциональный наполнитель, решая
одновременно экологическую и инженерную задачи {[}4{]}. Интерес к этому
подходу сформировался не вчера: ещё Pappu и соавт. (2007) показали
возможность замены традиционных наполнителей золой, древесной пылью и
красным шламом в полиэфирных композициях {[}5{]}. В дальнейшем это
направление развивалось в рамках исследований древесно-полимерных
композитов и наполнения полимеров минеральными порошками, а в 2018--2025
гг. наблюдается новый всплеск работ по ``зелёным'' композитам, что
хорошо видно по динамике публикаций. Современные работы охватывают
широкий спектр отходов: зола электростанций, доменный шлак, шлам
байеровского процесса, отходы текстиля, резиновых шин, печатных плат,
строительный бой, мраморная пыль и др. {[}1{]}. Также исследуются
различные полимерные матрицы - термопласты (ПЭ, ПП, ПВХ, полиамиды и
др.), реактопласты (эпоксидные, полиэфирные, фенольные смолы),
эластомеры (резины, силикон) - для понимания универсальности подхода.

Цель настоящего обзора систематизировать актуальные (2018-2025 гг.)
данные о применении техногенных отходов в роли наполнителей полимерных
композиционных материалов. Рассматриваются:

- масштабы проблемы и причины, по которым переработка отходов становится
технологической необходимостью;

- свойства наиболее распространённых техногенных наполнителей и подходы
к их подготовке и введению в полимерную матрицу;

- влияние состава и дисперсности наполнителя на ключевые характеристики
композитов;

- сопоставление результатов по разным типам отходов на основе единых
критериев; - примеры внедрённых решений и оценка их экологической и
экономической результативности.

Научная новизна обзора заключается в том, что он фокусируется на
публикациях последних лет и выполняет сопоставление разных групп
техногенных отходов в единой логике сравнения, что позволяет выделить
устойчивые тенденции и наиболее воспроизводимые практики.

{\bfseries Материалы и методы.} Настоящий обзор охватывает публикации
2018-2025 гг., включая работы в международных и отечественных журналах,
посвящённые полимерным композитам с техногенными наполнителями.
Приоритет отдан экспериментальным исследованиям, где приведены данные о
подготовке отходов (фракционирование, сушка, модификация),
характеристиках полимерных матриц и результатах испытаний механических,
физико-химических и термических свойств полученных композитов.
Дополнительно учтены публикации, обобщающие практический опыт применения
техногенного сырья в странах СНГ.

\emph{Полимерные матрицы.} В качестве связующей фазы в рассмотренных
работах используются как термопласты, так и термореактивные полимеры.

Среди термопластов наиболее часто встречаются полиолефины - полиэтилен
(ПЭ) и полипропилен (ПП), а также поливинилхлорид (ПВХ) {[}6{]},
полистирол (включая вспененный полистирол и отходы упаковки) и полиамиды
{[}7{]}. К термореактивным матрицам, применяемым в композитах с
техногенными наполнителями, преимущественно относятся эпоксидные,
ненасыщенные полиэфирные, фенолформальдегидные и полиуретановые смолы;
отдельный пласт работ посвящён эпоксидным системам с различными типами
отходов {[}8{]}. Для эластомерных матриц описаны силиконовые композиции
(например, системы «силикон-резина и красный шлам» для задач
ЭМИ-экранирования) {[}1{]}. В ряде исследований рассматриваются
материалы класса «полимерный бетон» и полимерцементные композиции на
основе полиэфирных или эпоксидных связующих с высокой долей минерального
наполнителя {[}8{]}.

\emph{Подготовка и модификация наполнителей.}

Перед введением в полимер техногенные отходы, как правило, подвергают
механической подготовке (дробление/помол до заданной дисперсности),
сушке и рассеву. Зола уноса ТЭС изначально представляет собой
тонкодисперсный порошок (обычно порядка 1-100 мкм) и во многих работах
используется без дополнительного помола; при необходимости выделяют
отдельные фракции, включая зольные микросфер полые алюмосиликатные
частицы сниженной плотности {[}9{]}. Доменный и сталеплавильный шлак,
напротив, чаще доводят помолом до порошка с размером частиц менее 100
мкм. Красный шлам, для которого характерна высокая доля частиц
<100 мкм, ддополнительно подвергают химической обработке --
промыванию и/или кислотной нейтрализации, чтобы снизить щёлочность и
повысить совместимость с полимерной матрицей. Стеклобой измельчают, в
том числе в шаровых мельницах до стеклянного порошка требуемой
дисперсности. Строительные отходы перерабатывают в мелкий заполнитель:
удаляют примеси и металл, выполняют дробление, а при необходимости --
тонкий помол, особенно для кирпичных отходов {[}10{]}. Для улучшения
адгезии к полимеру используют химическую модификацию поверхности:
обработку силонами, кислотами, нанесение наноструктурированных слоёв. В
термопластах широко применяют компатибилизаторы - сополимеры
(малеинированный ПП или ПЭ), которые при содержании 2-5\% улучшают
сцепление полярных минеральных и древесных наполнителей с матрицей. Для
эпоксидных смол используют предварительное пропитывание порошка
разбавленной смолой либо реагенты, связывающие гидроксильные группы
наполнителя с эпоксидом. Это позволяет уменьшить слабость межфазной
границы и предотвратить снижение прочности при росте доли наполнителя
{[}1{]}.

\emph{Методы изготовления композитов.}

Для термопластичных матриц наиболее распространены экструзия и литьё под
давлением. Порошкообразный наполнитель вводят в расплав полимера при
повышенной температуре, затем композицию гомогенизируют (часто с
использованием двухшнекового экструдера) {[}1{]}. После этого материал
гранулируют или формуют в виде профилей и листов. В лабораторной
практике образцы нередко получают методом горячего прессования из
предварительно смешанных компонентов. Для термореактивных систем
применяют заливку и пресс-формование. Наполнитель диспергируют в жидкой
смоле, вводят отвердитель, после чего смесь заливают в формы и проводят
отверждение при комнатной либо повышенной температуре, получая, как
правило, пластины для последующих испытаний {[}11{]}. В полимербетонах и
полимер-мортарных композициях доля минеральной фазы может достигать
70-80\%. Технология во многом аналогична приготовлению традиционных
строительных смесей, однако роль вяжущего здесь выполняет полимерная
композиция, отверждаемая воздушным или термическим способом.

Свойства композитов оценивают по стандартным методикам. Механические
характеристики обычно включают прочность при растяжении, изгибе и
сжатии, модуль упругости, ударную вязкость, твёрдость, износостойкость.
К физико-химическим показателям относят плотность, водопоглощение,
стойкость к растворителям, а также электроизоляционные и
электрофизические параметры.

Термическое поведение изучают через показатели теплостойкости и
температуры размягчения, а также методом термогравиметрического анализа
(TGA) для оценки термостабильности; для материалов строительного
назначения дополнительно анализируют горючесть и воспламеняемость
{[}2{]}. Полученные результаты сопоставляют с исходной полимерной
матрицей без наполнителя и с аналогичными композитами на традиционных
минеральных наполнителях, что позволяет корректно оценить вклад
техногенных добавок в изменение свойств полимерных материалов.

{\bfseries Обсуждение и результаты.} \emph{Виды техногенных наполнителей и
их влияние на свойства композитов.}

Анализ публикаций последних лет показывает, что техногенные отходы при
корректной подготовке и рациональном подборе концентрации могут
выступать эффективными наполнителями полимерных композитов. В ряде
случаев они обеспечивают рост отдельных эксплуатационных показателей, а
в остальных -- позволяют удерживать свойства на приемлемом уровне при
одновременном снижении материалоёмкости и себестоимости. Ниже
представлены основные группы техногенных наполнителей и обобщены
характерные эффекты их введения в различные полимерные матрицы; сводные
данные приведены в табл.1.
\end{multicols}

\tcap{Таблица 1 - Сравнительная характеристика техногенных отходов как наполнителей полимерных композитов}
\begin{longtblr}[
  label = none,
  entry = none,
  note{} = {\emph{{\bfseries Примечание:} стрелками ↑/↓ указано направление изменения
свойства при введении наполнителя относительно чистого полимера.
Конкретный эффект зависит от процентного содержания наполнителя и его
обработки; приведены типичные тенденции.}},
]{
  width = \linewidth,
  colspec = {Q[137]Q[217]Q[217]Q[365]},
  cells = {font = \small},
  row{1} = {c, font=\bfseries\small},
  hlines,
  vlines,
}
Вид техногенного наполнителя       & Основные характеристики                                     & Примеры полимерных матриц                                  & Влияние на свойства композита (М - механические, Ф - физико-химические, Т - термические)                                                                                  \\
Зола уноса (ТЭС)                   & Силикатно-алюминатная пыль; мелкие сферические частицы.     & ПЭ, ПП, ПВХ; эпоксидные и полиэфирные смолы; ДПК.          & {М: ↑ жесткость и прочность.\\Ф: ↓ усадка, слегка ↑ водопоглощение.\\Т: ↑ теплостойкость, ↓ горючесть}                                                                    \\
Доменный шлак                      & Кальций-алюмосиликатный порошок/зерна.                      & Полиамиды, ПВХ, полиэфирные смолы; полимербетон.           & {М: ↑ жесткость и износостойкость, возможное ↓ прочности на разрыв.Ф: возможное ↑ водопоглощения.\\Т: ↑ теплостойкость и огнестойкость}                                   \\
Красный шлам                       & Тонкий щелочной порошок оксидов Fe, Al и др.                & Эпоксидные, полиэфирные, винилэфирные смолы; спец. резины. & {М: ↑ твердость и прочность на сжатие.\\Ф: без обработки ухудшает влагостойкость, придаёт красно-коричневый цвет.Т: ↑ огнестойкость, ↓ дымообразование.}                  \\
Стеклобой                          & Аморфный SiO₂; остроугольные частицы.                       & Эпоксидные и полиэфирные композиты; ПП, ПЭ.                & {М: ↑ микротвердость и износостойкость; при больших дозах ↑ хрупкости.Ф: влагопоглощение ≈ 0; композит матовый.\\Т: ↑ теплостойкость, ↓ термическое расширение.}          \\
Строительный мусор (кирпич, бетон) & Пористый порошок дроблёных керамических и бетонных отходов. & Строительные полимерные композиты, ДПК; ПП, ПС.            & {М: ↑ модуль и прочность при сжатии; при избытке ↑ хрупкости.\\Ф: пористые фазы ↑ водопоглощение, керамика частично его снижает.Т: ↑ термостабильность и небольшое ↑ T₉.} 
\end{longtblr}

\begin{multicols}{2}
\emph{Зола уноса электростанций (пылеугольная зола).} Зола уноса -
мелкодисперсный минеральный порошок (1-100 мкм) на основе оксидов Si, Al
и Fe, с плотностью порядка 2,2-2,6 г/см³; частицы преимущественно
сферические, нередко присутствуют полые микросферы, что обеспечивает
высокую удельную поверхность и хорошую текучесть. Введение золы уноса в
полимерные матрицы исследовано наиболее широко. Механические свойства
для термопластов (ПЭ, ПП, ПВХ) добавка золы (до 5-15~мас. \%) обычно
повышает жесткость и твердость композита за счет заполнения полимерной
матрицы жесткими неорганическими частицами {[}12{]}.

Для композитов с зольным наполнителем в ряде работ фиксируется рост
прочности при растяжении и изгибе при умеренных дозировках, обычно в
пределах 5-10\% {[}13{]}. Так, добавление 6\% золы уноса в полиэфирную
матрицу сопровождалось увеличением прочности на разрыв на 65\% и
изгибной прочности на 32\%. Для полипропилена также описан рост ударной
вязкости: максимальный эффект достигался при около 9\% наполнения и
составлял порядка 74\% {[}13{]}. При дальнейшем повышении доли
наполнителя (более 20-30\%) прочностные характеристики нередко
ухудшаются, что связывают с агломерацией частиц и недостаточной адгезией
на границе раздела «частица-полимер» {[}1{]}. В термореактивных системах
зольный наполнитель, как правило, повышает модуль и твёрдость, однако
влияние на предельную прочность существенно зависит от качества
диспергирования и степени смачивания частиц смолой. Из физико-химических
эффектов важно учитывать гидрофильность золы: без поверхностной
модификации возможно некоторое увеличение водопоглощения, особенно при
наличии пористости. Одновременно зольные частицы способны снижать усадку
матрицы при отверждении или охлаждении за счёт заполнения объёма и
ограничения деформаций полимерной фазы {[}12{]}. Наличие тугоплавких
оксидов также может повышать термостабильность композитов и уменьшать
скорость терморазложения; в работе Shafik и соавт. (2025) показано, что
образцы с 40\% керамического порошка, близкого по природе к зольным
оксидным наполнителям теряли массу при нагреве медленнее и
демонстрировали несколько более высокую температуру начала разложения по
сравнению с чистым полистиролом {[}10{]}. Поскольку зола уноса
электрически непроводящая, её введение обычно не приводит к заметному
ухудшению электроизоляционных свойств.

\emph{Доменные и металлургические шлаки}. Доменный шлак
является побочным продуктом выплавки чугуна и, как правило, представлен
оксидной системой на основе CaO, SiO₂, Al₂O₃ и MgO при плотности порядка
2,8-3,0 г/см³. Для применения в полимерных композитах шлак обычно
измельчают до тонкости цемента. По характеру действия в полимерной
матрице шлаковый порошок во многом близок к золе уноса, поскольку также
состоит из твёрдых оксидных частиц. Наиболее устойчиво в литературе
отмечается рост прочности на сжатие и твёрдости, особенно в
полимербетонах и полиэфирных композициях строительного назначения
{[}14{]}. Например, для полиэфирного композита с 15\% шлака сообщается
увеличение прочности на сжатие на 4\%, а изгибной прочности на 21\%
{[}15{]}. В термопластичных матрицах эффекты зачастую более
неоднозначны: так, при введении гранулированного доменного шлака в
полиамид-6 на уровне 30-40\% наблюдали повышение износостойкости и
теплостойкости, однако прочность при растяжении снижалась, что объясняют
слабым межфазным сцеплением «шлак--полимер» при высоком наполнении
{[}1{]}. Шлаки, особенно базовые (CaO-rich), способны взаимодействовать
с влагой, поэтому композиты с ними могут проявлять умеренный рост
водопоглощения и в ряде случаев требуют добавок-стабилизаторов для
снижения возможных щелочных эффектов. Вместе с тем пористая структура
отдельных шлаковых частиц может способствовать поглощению энергии
деформации, что потенциально отражается на ударной вязкости. Термическая
стойкость композитов со шлаком обычно улучшается за счёт высокой
теплоёмкости и огнестойкости оксидного наполнителя. Дополнительно
подчёркивается, что доменные шлаки в большинстве случаев радиационно
безопасны и их применение в материалах не приводит к недопустимому росту
фонового излучения {[}2{]}.

\emph{Красный шлам}. Красный (бокситовый) шлам представляет
собой мелкодисперсный остаток производства глинозёма, содержащий оксиды
Fe₂O₃, Al₂O₃ и TiO₂, а также щелочные соединения; для него характерны
красно-бурый цвет и повышенная плотность порядка 3,3--3,6 г/см³. При
введении умеренных количеств красного шлама композиты обычно
демонстрируют рост твёрдости и модуля упругости. Так, в полиэфирной
матрице добавки 10-15\% шлама увеличивают твёрдость и прочность при
сжатии, одновременно повышая плотность материала {[}15{]}. Также
сообщается об улучшении износостойкости и сопротивления ползучести, что
связывают с вкладом твёрдых частиц Fe₂O₃. При чрезмерно высоком
наполнении (более 30\%) прочность при растяжении может снижаться,
поскольку частицы шлама нередко имеют пластинчатую морфологию и способны
выступать концентраторами напряжений при изгибе. Отдельный интерес
представляет влияние шлама на ударное поведение и демпфирование: в
исследовании {[}11{]} эпоксидно-бутадиеновая матрица с 50-90\% красного
шлама демонстрировала повышенную ударопрочность и поглощение энергии
удара по сравнению с композитом, наполненным глиноземистым порошком.
Щёлочность шлама в составе полимерной матрицы проявляется слабее, однако
при недостаточной промывке возможна деградация некоторых полимеров из-за
остаточной щёлочи. Композиты с красным шламом нередко уступают по
влагостойкости системам с более инертными наполнителями, поэтому могут
требовать гидрофобизаторов или иной поверхностной модификации. За счёт
высокого содержания Fe₂O₃ шлам способен повышать тепло- и огнестойкость;
сообщалось, что введение 20-30\% шлама в эпоксидные композиции уменьшает
скорость горения и дымообразование. Помимо этого, красный шлам придаёт
материалам функциональные свойства, включая электромагнитное
экранирование: показано, что смесь красного шлама с углеродными
нанотрубками в силиконовой матрице эффективно поглощает СВЧ-излучение и
может использоваться как экран электромагнитных помех {[}16{]}, что
иллюстрирует потенциал высокотехнологичных применений данного отхода.

\emph{Строительный мусор}. Отходы строительной индустрии
неоднородны по составу и включают фрагменты бетона, кирпича, керамики и
другие минеральные компоненты. Их плотность обычно находится в пределах
2,0--2,5 г/см³. После сортировки, удаления примесей и дробления материал
нередко дополнительно размалывают до порошка с размером частиц порядка
100--500 мкм; в таком виде он может применяться как минеральный
наполнитель полимерных композитов. По природе это сырьё близко к
традиционным наполнителям (песок, кварцевые порошки), поэтому его
введение обычно повышает модуль упругости и может умеренно увеличивать
прочность при сжатии. Вместе с тем прочность при растяжении и изгибе
иногда снижается из-за слабого межфазного контакта между полимером и
частицами «старого» цементного камня. Так, в работе Keskisaari и соавт.
(2016) добавление порошка кирпича и бетона в древесно-полимерный
композит привело к росту плотности и жёсткости, однако прочность в целом
сохранилась на уровне материала без наполнителя {[}17{]}. Более свежие
исследования показывают, что потенциал строительных отходов существенно
раскрывается при грамотном подборе фракционного состава и комбинации
компонентов. Shafik и соавт. (2025) получили композиты на основе отходов
упаковочного полистирола, усиливая их гибридными наполнителями из
древесной стружки, керамической крошки и кирпичного боя {[}10{]}.
Оптимальная композиция (20\% керамики и 20\% древесины) обеспечила более
чем двукратный рост прочности при растяжении, а керамическая
составляющая одновременно заметно снизила влагопоглощение и улучшила
термическую устойчивость: на термограммах образцы с кирпичным и
керамическим наполнителем характеризовались меньшей скоростью разложения
и более поздним началом деструкции {[}10{]}. В результате можно
заключить, что при корректной подготовке и подборе состава строительные
отходы, особенно керамическая фракция, способны выполнять роль
эффективного наполнителя для композитов строительного назначения,
повышая как прочностные, так и долговечностные характеристики. При этом
такие материалы сохраняют диэлектрические свойства, достаточные для
применения в изоляционных изделиях; в частности, для композитов с
кирпичным порошком отмечены низкие диэлектрические потери и проводимость
{[}10{]}.

\emph{Стеклобой (отходы стекла).} Битое стекло после помола представляет
собой стеклянный порошок, химически близкий к силикатным наполнителям,
примерно 70-75\% SiO₂ с присутствием Na₂O и CaO при плотности около 2,5
г/см³. По составу стеклопорошок сопоставим с кварцевыми наполнителями,
однако отличается гладкой поверхностью частиц, что может снижать
межфазное сцепление без специальных приёмов модификации. В композитах
стеклобой чаще всего ведёт себя как инертный минеральный наполнитель,
повышая твёрдость и улучшая стойкость к царапанию и истиранию. Например,
для эпоксидных покрытий показано, что введение тонкоизмельчённого стекла
увеличивает микротвёрдость по Виккерсу и повышает абразивную стойкость
по сравнению с чистой эпоксидной смолой {[}18{]}. Прочность при изгибе и
растяжении при умеренном содержании стеклопорошка (порядка 5-10\%) может
незначительно возрастать, тогда как при высоких долях (более 30\%)
возможно снижение прочности из-за концентрации напряжений у острых
граней стеклянных частиц {[}19{]}. Важным преимуществом стеклянного
наполнителя является крайне низкое влагопоглощение, поэтому композиты
обычно сохраняют невысокое водопоглощение; стекло также является
диэлектриком, и электрические изоляционные характеристики, как правило,
не ухудшаются. При этом оптические свойства пластика меняются ожидаемо:
прозрачность падает, материал становится мутным. Благодаря высокой
термостойкости стекло способно повышать теплостойкость композита и
сдвигать температурные характеристики в сторону более стабильного
поведения при нагреве (эффект зависит от матрицы и наполнения). Для
полимербетонов на полиэфирных или эпоксидных связующих замена части
кварцевого песка стеклобоем приводила к росту прочности при сжатии
(примерно на 5--10\%) и увеличению изгибной прочности (до +3 МПа)
{[}9{]}, что связывают с высокой прочностью стеклянных частиц и их
плотной упаковкой в матрице. Таким образом, стеклобой может
рассматриваться как практичная альтернатива коммерческим кварцевым
наполнителям и микрокальциту, переводя стекольные отходы в ресурс для
производства материалов.

\fig[\columnwidth]{c2/image8}[Рис.1 - Качественная оценка влияния
различных техногенных отходов на свойства полимерных композитов]

\emph{Прочие техногенные наполнители.} Помимо перечисленных, в последние
годы активно изучаются и другие потоки отходов. Резиновая крошка из
изношенных шин используется для повышения упругости и демпфирования и
находит применение в асфальтополимерных системах и ряде полимерных
композитов {[}20{]}. Текстильные отходы (например, измельчённые
джинсовые ткани) рассматривались как волокнистый наполнитель в
биоразлагаемых матрицах и в отдельных работах способствовали росту
прочности при растяжении {[}21{]}. Отходы электронной промышленности
также представляют интерес: после удаления металлов неметаллическая
фракция печатных плат измельчается и вводится в термопласты и
термореактивные смолы, повышая жёсткость и огнестойкость, хотя ударная
вязкость при этом может снижаться {[}22- 24{]}. Отдельно стоит отметить
применение золы и шлаков в геополимерных системах {[}23{]}; хотя это уже
относится к неорганическим матрицам, по сути, речь идёт о композитах,
где отходы замещают цементное сырьё. В рамках обзора, ориентированного
на полимерные матрицы, общий вывод остаётся одинаковым: большинство
минеральных порошковых отходов при соответствующей подготовке способны
работать как наполнители, обеспечивая свойства, сопоставимые с
традиционными минеральными добавками, а иногда и придавая материалам
дополнительные функциональные эффекты.

\emph{Сравнительные результаты физико-химических исследований полимерных
композитов с техногенными наполнителями.} Обобщение литературных данных
{[}1,5{]} показывает, что техногенные минеральные наполнители по-разному
влияют на ключевые физико-химические характеристики полимерных
композитов -- плотность, водопоглощение, усадку, термостабильность и
огнестойкость. В большинстве работ введение порошков отходов приводит к
росту плотности и уменьшению усадки по сравнению с чистой матрицей за
счёт замещения части объёма полимера более плотной неорганической фазой
и снижения доли усадочной составляющей {[}7, 8{]}.

По влиянию на водопоглощение и влагостойкость можно условно выделить две
группы систем. Относительно инертные наполнители -- зола уноса,
стеклобой при умеренном содержании (до \textasciitilde10-20 мас. \%)
лишь слабо увеличивают водопоглощение либо сохраняют его на уровне
исходного полимера благодаря низкой открытой пористости частиц и
практически отсутствующей собственной сорбционной способности {[}6,
9{]}. Наполнители с развитой пористостью и/или щелочной поверхностью,
такие как красный шлам и строительный мусор (бетонный и кирпичный бой),
напротив, вызывают заметный рост водопоглощения и снижение
влагостойкости композитов {[}3, 18{]}. Показано, что применение
поверхностных модификаторов и компатибилизаторов позволяет частично
нивелировать этот эффект и сохранить приемлемый уровень прочности
{[}20,21{]}. Большинство техногенных наполнителей улучшает
термостабильность и огнестойкость материалов. Зола уноса, шлаки,
стеклобой и красный шлам, состоящие преимущественно из тугоплавких
оксидов, повышают температуру начала термического разложения, уменьшают
скорость потери массы при нагреве и в ряде случаев снижают горючесть и
дымообразование {[}5, 9, 16{]}. Это особенно характерно для композитов
строительного и электротехнического назначения. При этом
электрофизические свойства многих систем (зола уноса, стеклобой,
значительная часть строительных отходов) остаются на уровне исходных
полимеров, что позволяет использовать такие композиты в качестве
конструкционно-изоляционных материалов {[}19, 21{]}.

Рассмотренные результаты подтверждают, что техногенные отходы могут
успешно выполнять функцию наполнителей, во многом аналогичную
традиционным (кальцит, тальк, стекловолокно), и в некоторых случаях
придавать композитам уникальные свойства. Встает вопрос, насколько эта
концепция реализована на практике и каковы ее преимущества с точки
зрения экологии и экономики?

С экологической точки зрения выгоды очевидны. Каждая тонна отходов,
вовлечённая во вторичное использование, это тонна, не отправленная на
полигон. Использование техногенных наполнителей снижает потребность в
добыче и производстве первичных наполнителей (карбонатных, силикатных),
а значит, уменьшает связанную с этим нагрузку на природу. Снижение
объема захороняемых отходов напрямую способствует защите почв и вод от
потенциального загрязнения тяжелыми металлами, которые присутствуют в
некоторых отходах. Заливка таких отходов в полимер фактически
иммобилизует опасные элементы в твердой матрице, предотвращая их
вымывание {[}24{]}. Например, композиты на основе красного шлама могут
служить барьером для тяжелых металлов вместо хранения жидкой шламовой
пульпы, ее можно включить в полимербетонные блоки, безопасные при
эксплуатации. В случае золы уноса использование в композитах устраняет
проблему ветровой эрозии золоотвалов. Производство полимеров обычно
сопровождается выбросами CO₂, но добавление наполнителя сокращает долю
полимера в продукте. ``Углеродный след'' отхода условно нулевой (или
даже отрицательный, учитывая, что в противном случае мог бы выделяться
метан при захоронении органических отходов). Таким образом, наполненные
отходами пластики могут иметь меньший углеродный след на единицу
изделия, чем чистые пластики.

Следует, впрочем, отметить, что экологический эффект должен
подтверждаться оценкой жизненного цикла. Не все отходы одинаково
безопасны. Например, зола и шлам могут содержать радиоактивные элементы.
Требуется контроль радиационной активности как упоминалось, многие золы
проходят по первому классу (А\tsb{эфф} <{} 370~Бк/кг)
и пригодны для использования {[}2{]}. Если же отход имеет повышенный
радиационный фон или сильно токсичен, его применение в массовых изделиях
ограничено санитарными нормами. Полимерные композиты с отходами должны
тестироваться на вымывание и эмиссию вредных веществ. В большинстве
случаев инкапсуляция в полимер надежно удерживает загрязнители.
Например, испытания геополимеров с красным шламом показали минимальное
выщелачивание тяжелых металлов {[}11{]}. Аналогично, зола уноса в
ПВХ-композите не высвобождает вредных веществ даже при контакте с водой,
что позволяет использовать такие материалы в несущих конструкциях и
несжатых изделиях без риска для окружающей среды.

Экономические выгоды двоякие. Для производителей материалов это
удешевление сырья, для предприятий-отходообразователей это снижение
расходов на утилизацию. Минеральные наполнители традиционно стоят
значительно дешевле полимеров. В случае техногенных отходов цена
наполнителя может быть вовсе нулевой или отрицательной, когда полигон
берет плату за прием отхода. Как отмечают исследователи, добавление
дешевого наполнителя вместо части полимера - это прямой путь к
удешевлению продукции {[}12{]}. Например, если 30\% дорогого эпоксидного
связующего заменить золой, стоимость композита снизится, а прочностные
характеристики могут даже вырасти. Это обусловлено тем, что
себестоимость заполненного композита ниже, а получаемые свойства все еще
находятся на приемлемом уровне. Более того, во многих случаях снижение
прочности вовсе не происходит, как мы видели, иногда прочность растет.
Таким образом, экономическая целесообразность доказана. Композиты на
отходах конкурируют по свойствам с обычными, но обходятся дешевле в
производстве. Для предприятия, генерирующего отход, создание продукта с
его использованием означает превращение пассива в актив. Например,
теплоэлектростанция, передавая золу в производство композитов, экономит
на эксплуатации золоотвалов и может получать доход от продажи золы как
сырья. В Алтайском крае уже практикуется передача золы ТЭЦ
дорожно-строительным предприятиям, что местные власти поддерживают как
более экологичный и выгодный способ обращения с отходом {[}25{]}.
Конечно, есть и затраты. Требуется оборудование для помола, сушки
отходов, добавки для улучшения совместимости, контроль качества. При
масштабном производстве они, как правило, окупаются. Экономический
анализ Барахтенко В.В. (2018) показал, что оптимизация рецептур ПВХ-зола
композита позволяет добиться прочности на уровне стандартов, при этом
снижение себестоимости достигает нескольких десятков процентов за счет
замещения ПВХ золой {[}2{]}.

Несмотря на очевидные плюсы, применение техногенных наполнителей требует
учета некоторых моментов. Во-первых, нестабильность состава отходов.
Качество, химический состав могут варьировать от партии к партии,
усложняя стандартизацию. В качестве решении - тщательная входная
характеристика сырья, возможно, смешение больших объемов для усреднения
параметров. Во-вторых, необходимость измельчения до тонкой фракции. Это
энергозатратно, и слишком крупный наполнитель ведет к плохим свойствам.
Однако современные мельницы и классификаторы способны готовить
ультрадисперсные порошки, более того, у многих отходов уже изначально
мелкая фракция. В-третьих, износ оборудования. Абразивные частицы, как
шлаки и стекла могут повышать износ экструдеров и смесителей. Здесь
пригодятся твердые наплавки на шнеки или выбор оборудования,
рассчитанного на минералонаполненные смеси. В-четвертых, логистика: для
эффективного внедрения желательно, чтобы поблизости от места образования
отхода было производство пластмасс или стройматериалов. Именно поэтому
особенно перспективна тема на стыке отраслей. Например, при ТЭС
организовать производство блоков или панелей из золонаполненного
композита, или при металлоплавильном заводе -- цех композитных изделий с
шлаком. Такой подход укладывается в концепцию промышленных симбиозов и
циркулярной экономики.

Как следует из обобщённых сравнительных данных по физико-химическим
характеристикам композитов с техногенными наполнителями {[}1-9,
13-23{]}, большинство таких систем обеспечивает повышение термостойкости
и жёсткости при допустимом росте водопоглощения.

Научные исследования продолжаются, разрабатываются нанокомпозиты с
отходами, комбинированные наполнители, а также биополимерные композиты с
техногенными отходами, что позволяет получать более экологичные продукты
на всём жизненном цикле {[}26{]}. Интерес представляет и аддитивное
производство: эксперименты с 3D-печатью пластиками, наполненными
порошками золы или древесных отходов, показали возможность создания
сложных деталей из таких материалов {[}27{]}.

{\bfseries Выводы.} Анализ работ 2018-2025 годов показывает, что
использование техногенных отходов в качестве наполнителей для полимерных
композитов -- это устойчивое и уже хорошо обоснованное направление,
позволяющее одновременно решать задачи утилизации отходов и снижения
расхода первичных ресурсов. Практически все основные типы полимерных
матриц могут быть успешно наполнены золой уноса, шлаками, красным
шламом, строительным мусором, стеклобоем и другими минеральными
порошками при условии правильной подготовки и модификации наполнителя. В
большинстве случаев техногенные наполнители повышают жёсткость, модуль
упругости, твердость и термостойкость композитов. Наиболее перспективной
областью применения остаются строительные и электротехнические
материалы. Экологический и экономический эффект связан с уменьшением
объёмов отходов, снижением нагрузки на полигоны и удешевлением продукции
за счёт замещения части полимера дешёвым вторичным сырьём. Дальнейшее
развитие направления связано со стандартизацией вторичных наполнителей,
накоплением инженерного опыта, применением нано- и поверхностных
модификаций для улучшения межфазного взаимодействия, а также с оценкой
жизненного цикла материалов.
\end{multicols}

\begin{center}
{\bfseries Литература}
\end{center}

\begin{refs}
1. Muthukutti G.P., Singh M.K., Palaniappan S.K., et al. Value-added
polymer composites using non-metallic industrial waste: A concise
review // Chemical Engineering Journal. -2025. -Vol.512: 162344. DOI
\href{https://ui.adsabs.harvard.edu/link_gateway/2025ChEnJ.51262344M/doi:10.1016/j.cej.2025.162344}{10.1016/j.cej.2025.162344}.

2. Барахтенко В.В., Бурдонов А.Е., Зелинская Е.В. Оценка эффективности
применения промышленных отходов в качестве наполнителя
поливинилхлоридных композиций // Строительство и реконструкция. -2018.
-№ 6 (80). - С.74--80.

3. Banda M. F., Matabane L., Munyengabe A. A phytoremediation approach
for the restoration of coal fly ash polluted sites: A review //
Heliyon. -2024. -Vol.10(23): e40741. DOI
\href{https://doi.org/10.1016/j.heliyon.2024.e40741}{10.1016/j.heliyon.2024.e40741}.

4. Барахтенко В.В. Строительный композиционный материал на основе отходов
ПВХ и золы уноса ТЭЦ: автореф. дис. по ВАК РФ 05.23.05.
Санкт-Петербург: СПбГАСУ. -2014.

5. Pappu A., Saxena M., Asolekar S. R. Solid wastes generation in India
and their recycling potential in building materials // Building and
Environment. -2007. -Vol.42(6). -P.2311-2320. DOI
\href{https://doi.org/10.1016/j.buildenv.2006.04.015}{10.1016/j.buildenv.2006.04.015}.

6. Evans L., Harrison P., Lawrence J., Casell J., Thiemann D., Kim J.
Coal Ash Primer/Earthjustice. -2023. -- 34 p.

7. Jagadeesh P., Puttegowda M., Thyavihalli Girijappa Y.G., Rangappa
S.M., Siengchin S. Effect of natural filler materials on fiber
reinforced hybrid polymer composites: An overview // Journal of
Natural Fibers. -2020. -Vol.19(11). -P.4132--4147. DOI
\href{https://doi.org/10.1080/15440478.2020.1854145}{10.1080/15440478.2020.1854145}.

8. Ağcan A.E., Kartal İ. A Review of Waste-Derived Fillers for Enhancing
the Properties of Epoxy Resins // International Journal of Adhesion
and Adhesives. -2025. -Vol.138: 103944. DOI
\href{https://doi.org/10.1016/j.ijadhadh.2025.103944}{10.1016/j.ijadhadh.2025.103944}.

9. Saribiyik M., Piskin A., Saribiyik A. The effects of waste glass
powder usage on polymer concrete properties // Construction and
Building Materials. -2013. -Vol.47. -P.840-844. DOI
\href{https://doi.org/10.1016/j.conbuildmat.2013.05.023}{10.1016/j.conbuildmat.2013.05.023}.

10. Shafik E.S., Wissa D.A., Labeeb A.M., Ward A.A., Abd-El-Messieh S.L.
Green polymeric composites based on construction and packaging waste:
toward advanced insulating materials // Scientific Reports. -2025.
-Vol.15:38495. DOI 10.1038/s41598-025-22450-z.

11. Melnyk L., Sviderskyy V. Development of multifunctional polymer
composites with high red mud content // Eastern-European Journal of
Enterprise Technologies. -2024. -Vol.6. -No.12(132). - P.34-43. DOI
\href{https://doi.org/10.15587/1729-4061.2024.317952}{10.15587/1729-4061.2024.317952}.

12. Colmar J., McBay A. Fly ash filler and polyvinyl chloride compositions
and conduits therefrom/ Patent publication number:~20020040084. -2002.

13. Nagaraja S., Anand P.B., Ammarullah M.I. Influence of fly ash filler
on the mechanical properties and water absorption behaviour of epoxy
polymer composites reinforced with pineapple leaf fibre for biomedical
applications // RSC Advances. -2024. -Vol.14(21). -P.14680-14696.
DOI:\href{https://doi.org/10.1039/d4ra00529e}{10.1039/d4ra00529e}.

14. Rachchh N.V., Misra R.K., Roychowdhary D.G. Effect of red mud filler
on mechanical and buckling characteristics of coir fibre-reinforced
polymer composite // Iranian Polymer Journal. -2015. -Vol.24(3). P.
253-265. DOI
\href{https://doi.org/10.1007/s13726-015-0317-4?urlappend=\%3Futm_source\%3Dresearchgate.net\%26utm_medium\%3Darticle}{10.1007/s13726-015-0317-4}.

15. Sarde B., Patil Y.D., Dholakiya B., Pawar V. Evaluation of effect of
red mud and fly ash as filler on the properties of modified PET resin
polymer mortar composite // Innovative Infrastructure Solutions.
-2025. -Vol.10(3). -P.1-12.
DOI\href{https://ui.adsabs.harvard.edu/link_gateway/2025InnIS..10..117S/doi:10.1007/s41062-025-01922-1}{10.1007/s41062-025-01922-1}.

16. Oral I., Özmeral N., Ahmetli G., Guzel H. Elastic and acoustic
properties determination of epoxy/polystyrene/nanoclay/red mud waste
hybrid composites by ultrasonic method // Polymer Composites. -2025.
-Vol.46(1). -P.930--948. DOI
\href{https://doi.org/10.1002/pc.29224}{10.1002/pc.29224}.

17. Keskisaari A., Butylina S., Kärki T. Use of construction and
demolition wastes as mineral fillers in hybrid wood--polymer
composites // Journal of Applied Polymer Science. -2016. -Vol.
133(19). DOI
\href{https://doi.org/10.1002/app.43412?urlappend=\%3Futm_source\%3Dresearchgate.net\%26utm_medium\%3Darticle}{10.1002/app.43412},

18. Buaphuen P., Makcharoen W., Vittayakorn W. The enhancement of polymer
composite coating by using a waste glass powder as alternative
reinforcement // Ferroelectrics. -2022. -Vol.601(1). -P.151-163. DOI
\href{https://doi.org/10.1080/00150193.2022.2130771?urlappend=\%3Futm_source\%3Dresearchgate.net\%26utm_medium\%3Darticle}{10.1080/00150193.2022.2130771}.

19. Bal B.C., Narlioglu N. Effect of waste glass reinforcements particle
size on some properties of HDPE-based composites // Tribology and
Materials. -2025. -Vol.4(1). -P.29--37.
DOI:\href{https://doi.org/10.46793/tribomat.2025.001}{10.46793/tribomat.2025.001}.

20. Abdulhameed J.I., Ali A.H., Kara İ.H., Mahan H.M., Konovalov S.V.,
Al-Nedawi N.M. Preparing eco-friendly composite from end-life tires
and epoxy resin and examining its mechanical, and acoustic insulation
properties // International Journal of Nanoelectronics and Materials.
-2024. -Vol.17(1). -P.1-5. DOI
\href{https://doi.org/10.58915/ijneam.v17i1.514}{10.58915/ijneam.v17i1.514}.

21. Shafqat A.R., Hussain M., Nawab Y., Ashraf M., Ahmad S., Batool G.
Circularity in materials: a review on polymer composites made from
agriculture and textile waste // International Journal of Polymer
Science. -2023. -Vol.2023(1): 5872605. DOI
\href{https://doi.org/10.1155/2023/5872605}{10.1155/2023/5872605}.

22. Wang X., Fu C., Feng Z., Huo H., Yin X., Gao G., et al. Flyash/polymer
composite electrolyte with internal binding interaction enables
highly-stable extrinsic-interfaces of all-solid-state lithium
batteries // Chemical Engineering Journal. -2022. -Vol.428: 131041.
DOI
\href{https://doi.org/10.1016/j.cej.2021.131041}{10.1016/j.cej.2021.131041}.

23. Герасимова Н.П. Установка для интенсификации процесса затвердевания
бетона с использованием золы уноса ТЭЦ // Вестник ИрГТУ. -2016. -№
7(114). -C.126-132. DOI 10.21285/1814-3520-2016-7-126-132.

24. Flores-Campos R., Deaquino-Lara R., Rodríguez-Reyes M.,
Martínez-Sánchez R., Estrada-Ruiz R.H. The use of nonmetallic fraction
particles with the double purpose of increasing the mechanical
properties of low-density polyethylene composite and reducing the
pollution associated with the recycling of metals from e-waste //
Recycling. -2024. -Vol.9(4): 56. DOI
\href{https://doi.org/10.3390/recycling9040056}{10.3390/recycling9040056}.

25. Кузнецова Т.В. В Алтайском крае используют золу Барнаульской ТЭЦ для
ремонта дорог// Российская газета. -2025. - ЭЛ~№~ФС~77 -50379.

26. Chaturvedi A.K., Gupta M.K., Pappu A. The role of carbon nanotubes on
flexural strength and dielectric properties of water sustainable fly
ash polymer nanocomposites // Physica B: Condensed Matter. -2021.-Vol.
620: 413283. DOI 10.1016/j.physb.2021.413283.

27. Mojapelo K.S., Kupolati W.K., Burger E.A., Ndambuki J.M., Snyman J.,
Sadiku E.R. Recent progress in preparation, improvement and
applications of green polymer composites // Biopolymers. -2025. -Vol.
116(6): e70059. DOI
\href{https://doi.org/10.1002/bip.70059}{10.1002/bip.70059}.
\end{refs}

\begin{center}
{\bfseries References}
\end{center}

\begin{refs}
1. Muthukutti G.P., Singh M.K., Palaniappan S.K., et al. Value-added
polymer composites using non-metallic industrial waste: A concise review
// Chemical Engineering Journal. -2025. -Vol.512: 162344. DOI
\href{https://ui.adsabs.harvard.edu/link_gateway/2025ChEnJ.51262344M/doi:10.1016/j.cej.2025.162344}{10.1016/j.cej.2025.162344}.

2. Barahtenko V.V., Burdonov A.E., Zelinskaja E.V. Ocenka jeffektivnosti
primenenija promyshlennyh othodov v kachestve napolnitelja
polivinilhloridnyh kompozicij // Stroitel' stvo i
rekonstrukcija. -2018. -№ 6 (80). -S.74--80. {[}in Russian{]}

3. Banda M. F., Matabane L., Munyengabe A. A phytoremediation approach
for the restoration of coal fly ash polluted sites: A review //Heliyon.
-2024. -Vol.10(23): e40741. DOI 10.1016/j.heliyon.2024.e40741. {[}in
Russian{]}

4. Barahtenko V.V. Stroitel' nyj kompozicionnyj material
na osnove othodov PVH i zoly unosa TJeC: avtoref. dis. po VAK RF
05.23.05. Sankt-Peterburg: SPbGASU. -2014. {[}in Russian{]}

5. Pappu A., Saxena M., Asolekar S. R. Solid wastes generation in India
and their recycling potential in building materials //Building and
Environment. -2007. -Vol.42(6). -P.2311--2320. DOI
\href{https://doi.org/10.1016/j.buildenv.2006.04.015}{10.1016/j.buildenv.2006.04.015}.

6. Evans L., Harrison P., Lawrence J., Casell J., Thiemann D., Kim J.
Coal Ash Primer/Earthjustice. -2023. -34 p.

7. Jagadeesh P., Puttegowda M., Thyavihalli Girijappa Y.G., Rangappa
S.M., Siengchin S. Effect of natural filler materials on fiber
reinforced hybrid polymer composites: An overview // Journal of Natural
Fibers. -2020. -Vol.19(11). -P.4132--4147. DOI
\href{https://doi.org/10.1080/15440478.2020.1854145}{10.1080/15440478.2020.1854145}.

8. Ağcan A.E., Kartal İ. A Review of Waste-Derived Fillers for Enhancing
the Properties of Epoxy Resins // International Journal of Adhesion and
Adhesives. -2025. -Vol.138: 103944. DOI
\href{https://doi.org/10.1016/j.ijadhadh.2025.103944}{10.1016/j.ijadhadh.2025.103944}.

9. Saribiyik M., Piskin A., Saribiyik A. The effects of waste glass
powder usage on polymer concrete properties // Construction and Building
Materials. -2013. -Vol.47. -P.840-844. DOI
\href{https://doi.org/10.1016/j.conbuildmat.2013.05.023}{10.1016/j.conbuildmat.2013.05.023}.

10. Shafik E.S., Wissa D.A., Labeeb A.M., Ward A.A., Abd-El-Messieh S.L.
Green polymeric composites based on construction and packaging waste:
toward advanced insulating materials // Scientific Reports. -2025. -Vol.
15:38495. DOI 10.1038/s41598-025-22450-z.

11. Melnyk L., Sviderskyy V. Development of multifunctional polymer
composites with high red mud content // Eastern-European Journal of
Enterprise Technologies. -2024. -Vol.6. -No.12(132). -P.34-43. DOI
\href{https://doi.org/10.15587/1729-4061.2024.317952}{10.15587/1729-4061.2024.317952}.

12. Colmar J., McBay A. Fly ash filler and polyvinyl chloride
compositions and conduits therefrom/ Patent publication
number:~20020040084. -2002.

13. Nagaraja S., Anand P.B., Ammarullah M.I. Influence of fly ash filler
on the mechanical properties and water absorption behaviour of epoxy
polymer composites reinforced with pineapple leaf fibre for biomedical
applications // RSC Advances. -2024. -Vol.14(21). -P.14680-14696.
DOI\href{https://doi.org/10.1039/d4ra00529e}{10.1039/d4ra00529e}.

14. Rachchh N.V., Misra R.K., Roychowdhary D.G. Effect of red mud filler
on mechanical and buckling characteristics of coir fibre-reinforced
polymer composite // Iranian Polymer Journal. -2015. -Vol.24(3). P.
253-265. DOI
\href{https://doi.org/10.1007/s13726-015-0317-4?urlappend=\%3Futm_source\%3Dresearchgate.net\%26utm_medium\%3Darticle}{10.1007/s13726-015-0317-4}.

15. Sarde B., Patil Y.D., Dholakiya B., Pawar V. Evaluation of effect of
red mud and fly ash as filler on the properties of modified PET resin
polymer mortar composite // Innovative Infrastructure Solutions. -2025.
-Vol.10(3). -P.1--12. DOI
\href{https://ui.adsabs.harvard.edu/link_gateway/2025InnIS..10..117S/doi:10.1007/s41062-025-01922-1}{10.1007/s41062-025-01922-1}.

16. Oral I., Özmeral N., Ahmetli G., Guzel H. Elastic and acoustic
properties determination of epoxy/polystyrene/nanoclay/red mud waste
hybrid composites by ultrasonic method // Polymer Composites. -2025.
-Vol.46(1). - P.930-948. DOI
\href{https://doi.org/10.1002/pc.29224}{10.1002/pc.29224}.

17. Keskisaari A., Butylina S., Kärki T. Use of construction and
demolition wastes as mineral fillers in hybrid wood--polymer composites
// Journal of Applied Polymer Science. -2016. -Vol.133(19). DOI
\href{https://doi.org/10.1002/app.43412?urlappend=\%3Futm_source\%3Dresearchgate.net\%26utm_medium\%3Darticle}{10.1002/app.43412},

18. Buaphuen P., Makcharoen W., Vittayakorn W. The enhancement of
polymer composite coating by using a waste glass powder as alternative
reinforcement // Ferroelectrics. -2022. -Vol.601(1). -P.151-163. DOI
\href{https://doi.org/10.1080/00150193.2022.2130771?urlappend=\%3Futm_source\%3Dresearchgate.net\%26utm_medium\%3Darticle}{10.1080/00150193.2022.2130771}.

19. Bal B.C., Narlioglu N. Effect of waste glass reinforcements particle
size on some properties of HDPE-based composites // Tribology and
Materials. -2025. -Vol.4(1). -P.29-37. DOI
\href{https://doi.org/10.46793/tribomat.2025.001}{10.46793/tribomat.2025.001}.

20. Abdulhameed J.I., Ali A.H., Kara İ.H., Mahan H.M., Konovalov S.V.,
Al-Nedawi N.M. Preparing eco-friendly composite from end-life tires and
epoxy resin and examining its mechanical, and acoustic insulation
properties // International Journal of Nanoelectronics and Materials.
-2024. -Vol.17(1). -P.1-5. DOI
\href{https://doi.org/10.58915/ijneam.v17i1.514}{10.58915/ijneam.v17i1.514}.

21. Shafqat A.R., Hussain M., Nawab Y., Ashraf M., Ahmad S., Batool G.
Circularity in materials: a review on polymer composites made from
agriculture and textile waste // International Journal of Polymer
Science. -2023. -Vol.2023(1):5872605. DOI
\href{https://doi.org/10.1155/2023/5872605}{10.1155/2023/5872605}.

22. Wang X., Fu C., Feng Z., Huo H., Yin X., Gao G., et al.
Flyash/polymer composite electrolyte with internal binding interaction
enables highly-stable extrinsic-interfaces of all-solid-state lithium
batteries // Chemical Engineering Journal. -2022. -Vol.428: 131041. DOI
\href{https://doi.org/10.1016/j.cej.2021.131041}{10.1016/j.cej.2021.131041}.

23. Gerasimova N.P. Ustanovka dlja intensifikacii processa
zatverdevanija betona s ispol' zovaniem zoly unosa TJeC
// Vestnik IrGTU. -2016. -№ 7 (114). --S.126-132. DOI
10.21285/1814-3520-2016-7-126-132. {[}in Russian{]}

24. Flores-Campos R., Deaquino-Lara R., Rodríguez-Reyes M.,
Martínez-Sánchez R., Estrada-Ruiz R.H. The use of nonmetallic fraction
particles with the double purpose of increasing the mechanical
properties of low-density polyethylene composite and reducing the
pollution associated with the recycling of metals from e-waste //
Recycling. -2024. -Vol.9(4): 56. DOI
\href{https://doi.org/10.3390/recycling9040056}{10.3390/recycling9040056}.

25. Kuznecova T.V. V Altajskom krae ispol' zujut zolu
Barnaul' skoj TJeC dlja remonta dorog// Rossijskaja
gazeta. -2025. - JeL № FS 77 -50379. {[}in Russian{]}

26. Chaturvedi A.K., Gupta M.K., Pappu A. The role of carbon nanotubes
on flexural strength and dielectric properties of water sustainable fly
ash polymer nanocomposites // Physica B: Condensed Matter. -2021.-Vol.
620: 413283. DOI 10.1016/j.physb.2021.413283.

27. Mojapelo K.S., Kupolati W.K., Burger E.A., Ndambuki J.M., Snyman J.,
Sadiku E.R. Recent progress in preparation, improvement and applications
of green polymer composites // Biopolymers. -2025. -Vol.116(6): e70059.
DOI \href{https://doi.org/10.1002/bip.70059}{10.1002/bip.70059}.
\end{refs}

\begin{info}
\hspace{1em}\emph{{\bfseries Сведения об авторах}}

Аманғалиева А.Е. - магистрант, Актюбинский региональный университет
им.  К. Жубанова, Актобе, Казахстан, e-mail: Ashlen002@gmail.com;

Орынбасар Р.О. - кандидат химических наук, доцент, Актюбинский
региональный университет им. К. Жубанова, Актобе, Казахстан, e-mail:
raihan\_06\_79@mail.ru;

Тастанова Л.К. - кандидат химических наук, профессор Heriot Watt
Aktobe Campus, Актюбинский региональный университет им. К. Жубанова,
Актобе, Казахстан, e-mail: lyazzatt@mail.ru;

Жакупова Г.Ж. - преподаватель, Актюбинский региональный университет
им.  К. Жубанова, Актобе, Казахстан, e-mail:
gulmira.zhakupova80@bk.ru;

~Жумабекова А.К. - кандидат химических наук, доцент, Казахский
университет технологии и бизнеса им.К.Кулажанова, Астана, Казахстан,
e-mail:,zhumabekova\_ak@mail.ru.

\hspace{1em}\emph{{\bfseries Information about the authors}}

Amangalieva A.Y. - Master' s student, K.Zhubanov Aktobe Regional
University, Aktobe, Kazakhstan, e-mail: Ashlen002@gmail.com;

Orynbasar R.O. - Candidate of Chemical Sciences, Associate Professor,
K.Zhubanov Aktobe Regional University, Aktobe, Kazakhstan, e-mail:
raihan\_06\_79@mail.ru;

Tastanova L.K. - Candidate of Chemical Sciences, Professor of Heriot
Watt Aktobe Campus, K.Zhubanov Aktobe Regional University, Aktobe,
Kazakhstan, e-mail: lyazzatt@mail.ru;

Zhakupova G.Z. - Lecturer K.Zhubanov Aktobe Regional University,
Aktobe, Kazakhstan, e-mail: gulmira.zhakupova80@bk.ru;

Zhumabekova A.K. - Candidate of Chemical Sciences, Associate
Professor, K.Kulazhanov Kazakh University of Technology and Business,
Astana, Kazakhstan, e-mail: zhumabekova\_ak@mail.ru.
\end{info}
